\chapterhead{42}{ORR-3}{Risk Identification}{1}{2}

\lohead{42}{a}{Compare different top-down and bottom-up approaches and tools for
identifying operational risks.}

\orsub{Operational risk identification with two approaches}

\begin{center}
\begin{tikzpicture}[font=\fontsize{8.4}{10.5}\selectfont,
  ptitle/.style={text width=6.9cm,align=center,rounded corners=3pt,inner sep=6pt,
                 text=white,font=\sffamily\bfseries\fontsize{9}{11}\selectfont},
  pbody/.style={text width=6.9cm,align=left,rounded corners=3pt,inner sep=8pt,
                fill=paleblue,draw=palenavy,line width=0.6pt,text=navy}]

\node[ptitle,fill=navy,anchor=north] (t1) at (-4.1,0) {1. Top-Down Approach};
\node[ptitle,fill=forest,anchor=north] (t2) at (4.1,0) {2. Bottom-Up Approach};

\node[pbody,anchor=north] (b1) at (t1.south) {
  Starts with senior management considering the bank's overall strategy and
  potential threats. Involves brainstorming sessions with key personnel from
  various departments. Identifies significant risks that could impact the bank's
  business goals.};

\node[pbody,anchor=north] (b2) at (t2.south) {
  Involves employees at lower and middle management levels. It is often conducted
  during Risk and Control Self-Assessment (RCSA); it focuses on individual
  business unit risks or processes within the bank.};

\draw[-{Stealth[length=5pt]},navy,line width=1.4pt]
  ($(b1.south west)+(0.5,-0.15)$) -- ($(b1.south east)+(-0.5,-0.15)$)
  node[midway,below,font=\sffamily\fontsize{7.2}{9}\selectfont,text=navy]
  {strategy $\rightarrow$ business units};
\draw[-{Stealth[length=5pt]},forest,line width=1.4pt]
  ($(b2.south east)+(-0.5,-0.15)$) -- ($(b2.south west)+(0.5,-0.15)$)
  node[midway,below,font=\sffamily\fontsize{7.2}{9}\selectfont,text=forest]
  {processes $\rightarrow$ firmwide view};
\end{tikzpicture}
\end{center}
\orfigcap{The two directions from which operational risk is identified.}

\orsub{Top-down risk identification tools}

\begin{enumerate}
\item[\textcolor{rust}{\sffamily\bfseries 1)}] \textbf{Exposures.} Exposures in
risk identification refer to critical areas like major clients, revenue sources,
and third-party dependencies that can lead to operational risks if they fail.
\item[\textcolor{rust}{\sffamily\bfseries 2)}] \textbf{Vulnerabilities.}
Weaknesses within the bank, like outdated processes, control deficiencies, or
neglected areas, can amplify the impact of an exposure if exploited.
\item[\textcolor{rust}{\sffamily\bfseries 3)}] \textbf{Risk wheel.} A visual tool
depicting various risk categories and their interconnections. Example: a natural
disaster (event) can lead to service outages (business continuity), impacting the
bank's reputation. It helps understand root causes of risks and prioritise them
for mitigation.
\item[\textcolor{rust}{\sffamily\bfseries 4)}] \textbf{Emerging risks.} Emerging
risks are a growing concern due to market volatility, technology advancements,
and cybersecurity issues. Two categories: predictable (cyber, regulatory) and
unpredictable (global pandemic). Horizon scanning (PESTLE analysis) helps
identify future threats. Emerging risk committees monitor trends and conduct
scenario analysis.
\end{enumerate}

\orsub{Bottom-up risk identification tools}

\orsubb{1) Event and loss data analysis}

Utilises historical data to forecast potential future losses and assess
operational risk concentrations.

\begin{orkeybox}
{\sffamily\bfseries\fontsize{8.6}{10}\selectfont\color{navy}Use of data to
analyse at bottom level}\par
\vspace{4pt}
\begin{lettered}
\item \textbf{Internal losses:} identify areas with high risk concentration
(back office, transactions).
\item \textbf{External losses} (from ORX data): provide insights into potential
risks and necessary control adjustments.
\item \textbf{Near misses:} highlight control gaps without incurred costs,
serving as valuable learning opportunities.
\end{lettered}
\end{orkeybox}

\orsubb{2) Risk and Control Self-Assessment (RCSA)}

\begin{lettered}
\item A self-evaluation tool for identifying key operational risks, evaluating
existing controls, and identifying residual risks not addressed by current
controls.
\item To do RCSA, involve experienced and new staff for diverse perspectives.
\item Often performed annually, or after major changes in risk profile or
external environment.
\item Avoid excessive detail to maintain focus on valuable insights.
\end{lettered}

\orsubb{3) Process mapping}

\begin{lettered}
\item Identify potential problems within each step of a business process.
\item Analyse how controls mitigate risks at each step.
\item Ensure the level of detail provides useful risk insights.
\item Aim for balance: each step linked to a key action and its control(s).
\end{lettered}

% ---------------------------------------------------------------------
\lohead{42}{b}{Describe best practices in the process of scenario analysis for
operational risk.}

\orsub{Best practices in scenario analysis for operational risk}

\begin{center}
\begin{tikzpicture}[font=\fontsize{8.2}{10.4}\selectfont,
  qt/.style={text width=7.4cm,align=left,rounded corners=3pt,inner sep=5pt,
             text=white,font=\sffamily\bfseries\fontsize{8.6}{10.5}\selectfont},
  qb/.style={text width=7.4cm,align=left,rounded corners=3pt,inner sep=7pt,
             draw=palenavy,line width=0.6pt,text=navy}]

\node[qt,fill=navy,anchor=north west] (a1) at (-8.15,0) {Regulatory guidance};
\node[qb,fill=paleblue,anchor=north west] (a2) at (-8.15,-0.62) {
  a)~Focuses on ``low probability, high severity'' events.\\[2pt]
  b)~Recommends business continuity plans (BCPs) with scenario analysis.\\[2pt]
  c)~Considers unforeseen incidents that could disrupt business.\\[2pt]
  d)~BCPs should consider:
  \begin{itemize}[leftmargin=12pt,itemsep=1pt,topsep=2pt,
        label=\textcolor{forest}{\scriptsize$\checkmark$}]
    \item Key business operations
    \item Reliance on internal/external parties (e.g., IT providers)
  \end{itemize}};

\node[qt,fill=teal,anchor=north west] (b1) at (0.35,0) {Brainstorming methods};
\node[qb,fill=paleblue,anchor=north west] (b2) at (0.35,-0.62) {
  a)~Aim for neutrality to avoid bias.\\[2pt]
  b)~Use evidence and detailed documentation to support scenarios.\\[2pt]
  c)~Consider providing participants with documents before a session:
  \begin{itemize}[leftmargin=12pt,itemsep=1pt,topsep=2pt,
        label=\textcolor{forest}{\scriptsize$\checkmark$}]
    \item Past loss data (internal and external)
    \item RCSA feedback
    \item Key risk indicator scores
  \end{itemize}
  d)~Participants:
  \begin{itemize}[leftmargin=12pt,itemsep=1pt,topsep=2pt,
        label=\textcolor{forest}{\scriptsize$\checkmark$}]
    \item Upper management from various business units
    \item Ideally, external risk specialists
  \end{itemize}};

\begin{scope}[on background layer]
  \node[fit=(a2)(b2),inner sep=0pt] (rowone) {};
\end{scope}

\node[qt,fill=forest,anchor=north west] (c1) at ($(rowone.south west)+(0,-0.3)$)
  {Scenario generation};
\node[qb,fill=palegreen,anchor=north west] (c2) at ($(c1.south west)+(0,0)$) {
  a)~Brainstorm many potential scenarios.\\[2pt]
  b)~Use facilitators.\\[2pt]
  c)~Consider past losses and potential future events.\\[2pt]
  d)~Use silent voting to avoid bias from dominant personalities.};

\node[qt,fill=olive,anchor=north west] (d1) at ($(rowone.south east)+(-7.5,-0.3)$)
  {Scenario selection};
\node[qb,fill=palegreen,anchor=north west] (d2) at ($(d1.south west)+(0,0)$) {
  a)~Classify scenarios by risk type or impact.\\[2pt]
  b)~Combine scenarios with similar impacts from different causes.\\[2pt]
  c)~Drop scenarios with immaterial impact.\\[2pt]
  d)~Select a final list for further assessment (e.g., 15--30 scenarios).};
\end{tikzpicture}
\end{center}
\orfigcap{The four practice areas of operational risk scenario analysis.}

% ---------------------------------------------------------------------
\lohead{42}{c}{Describe and apply an operational risk taxonomy and give examples
of different taxonomies of operational risks.}

\lohead{42}{d}{Describe and apply the Level 1, 2, and 3 categories in the Basel
operational risk taxonomy.}

\orsub{Operational risk taxonomies}

\orsubb{Basel taxonomy}

The Basel Committee on Banking Supervision (BCBS) has established an operational
risk taxonomy for banks, consisting of three levels.

\begin{center}
\begin{tikzpicture}[font=\fontsize{8.4}{10.5}\selectfont]
\node[fill=navy,text=white,rounded corners=3pt,inner sep=7pt,text width=3.1cm,
      align=center,font=\sffamily\bfseries\fontsize{9}{11}\selectfont] (L1) at (0,0)
      {Level 1\\[1pt]\mdseries\fontsize{7.6}{9}\selectfont 7 event types};
\node[fill=teal,text=white,rounded corners=3pt,inner sep=7pt,text width=5.0cm,
      align=center,font=\sffamily\bfseries\fontsize{9}{11}\selectfont] (L2) at (0,-1.95)
      {Level 2\\[1pt]\mdseries\fontsize{7.6}{9}\selectfont 20 subcategories};
\node[fill=forest,text=white,rounded corners=3pt,inner sep=7pt,text width=6.9cm,
      align=center,font=\sffamily\bfseries\fontsize{9}{11}\selectfont] (L3) at (0,-3.90)
      {Level 3\\[1pt]\mdseries\fontsize{7.6}{9}\selectfont detailed examples};

\node[text width=7.6cm,align=left,anchor=west,text=navy] at (4.1,0)
  {Provides a broad description of seven operational risk events/risks.};
\node[text width=7.6cm,align=left,anchor=west,text=navy] at (4.1,-1.95)
  {Offers greater specificity by breaking down risks into 20 specific risk
   subcategories. \textbf{Regulated banks are required to report operational risk
   using Level 1 and Level 2 categories.}};
\node[text width=7.6cm,align=left,anchor=west,text=navy] at (4.1,-3.90)
  {Provides even more detailed examples of risks but is not considered a
   regulatory requirement. It is used for analytical and logistical purposes.};
\end{tikzpicture}
\end{center}
\orfigcap{The three levels of the Basel operational risk taxonomy.}

\orsubb{Basel II operational risk categories and subcategories}

\noindent\begin{minipage}{\textwidth}
\renewcommand{\arraystretch}{1.3}
\begin{tabularx}{\textwidth}{@{}L{0.28\textwidth}L{0.30\textwidth}Y@{}}
\rowcolor{navy} \thd{Category} & \thd{Subcategory} & \thd{Example}\\
{\tabfont \textbf{1) Execution, Delivery, and Process Management (EDPM)}} &
{\tabfont Transaction capture, execution and maintenance} &
{\tabfont Data errors, miscommunications, delivery failure}\\
\rowcolor{paleblue}
{\tabfont \textbf{2) Clients, Products, and Business Practices}} &
{\tabfont Suitability, disclosure and fiduciary duties} &
{\tabfont Disclosure issues, privacy violations}\\
{\tabfont \textbf{3) Business Disruption and System Failures (BDSF)}} &
{\tabfont Hardware, software, telecommunications, and utility outage} &
{\tabfont Funding system has crashed, requiring \$30 billion in additional
 funding due to higher costs}\\
\rowcolor{paleblue}
{\tabfont \textbf{4) Internal Fraud}} &
{\tabfont Unauthorised activity} &
{\tabfont Not reporting transactions intentionally}\\
{\tabfont \textbf{5) External Fraud}} &
{\tabfont Theft and fraud} &
{\tabfont Theft, forgery, check kiting}\\
\rowcolor{paleblue}
{\tabfont \textbf{6) Employment Practices and Workplace Safety (EPWS)}} &
{\tabfont Employee relations} &
{\tabfont Compensation, benefits, termination, organised labour}\\
{\tabfont \textbf{7) Damage to Physical Assets (DPA)}} &
{\tabfont Disasters and other events} &
{\tabfont Natural disasters, human losses from external sources}\\
\end{tabularx}
\end{minipage}
\orfigcap{Level 1 categories with a Level 2 subcategory and Level 3 example each.}

\orsub{Taxonomies for causes, impacts, and controls}

\begin{center}
\begin{tikzpicture}[font=\fontsize{8.2}{10.4}\selectfont,
  ht/.style={text width=7.4cm,align=left,rounded corners=3pt,inner sep=5pt,
             text=white,font=\sffamily\bfseries\fontsize{8.6}{10.5}\selectfont},
  hb/.style={text width=7.4cm,align=left,rounded corners=3pt,inner sep=7pt,
             draw=palenavy,line width=0.6pt,text=navy,fill=paleblue}]

\node[ht,fill=rust,anchor=north west] (c1) at (-8.15,0) {Causes};
\node[hb,anchor=north west] (c2) at (-8.15,-0.62) {
  Going back to the Basel definition of operational risk provides the basic
  taxonomy for causes: employees, process failure, external, and systems.\\[4pt]
  $\blacktriangleright$~\textbf{Employees} --- human error, unintentional harm;
  incompetence and/or poor training\\[2pt]
  $\blacktriangleright$~\textbf{Process failure} --- failures in system design,
  process design, or governance\\[2pt]
  $\blacktriangleright$~\textbf{External} --- political, regulatory, economic,
  and physical aspects of the business environment\\[2pt]
  $\blacktriangleright$~\textbf{Systems} --- malfunctions, lack of capacity;
  insufficient maintenance and/or testing};

\node[ht,fill=teal,anchor=north west] (i1) at (0.35,0) {Impacts};
\node[hb,anchor=north west] (i2) at (0.35,-0.62) {
  Impacts are classified at Level 1 as direct financial (losses or remediation),
  indirect financial, and nonfinancial. Examples of the impacts at Level 2 are
  provided underneath each of the impacts at Level 1.\\[4pt]
  $\blacktriangleright$~\textbf{Direct financial impact} --- revenue reductions;
  lost or damaged assets; legal costs, fines and penalties\\[2pt]
  $\blacktriangleright$~\textbf{Indirect financial impact} --- regulatory
  compliance; reputational harm\\[2pt]
  $\blacktriangleright$~\textbf{Nonfinancial impact} --- customers, employees,
  third parties (e.g., suppliers, service providers), shareholders, competitors};

\begin{scope}[on background layer]
  \node[fit=(c2)(i2),inner sep=0pt] (toprow) {};
\end{scope}

\node[ht,fill=forest,anchor=north] (k1) at ($(toprow.south)+(0,-0.75)$) {Controls};
\node[hb,fill=palegreen,anchor=north] (k2) at (k1.south) {
  $\blacktriangleright$~\textbf{Preventive controls:} proactive measures to
  prevent risks. \emph{Example: edit checks in data entry to reduce errors.}\\[3pt]
  $\blacktriangleright$~\textbf{Detective controls:} used during or after an
  event to minimise negative impacts. \emph{Example: periodic transaction
  reconciliation to ensure accuracy.}\\[3pt]
  $\blacktriangleright$~\textbf{Corrective controls:} aim to minimise and rectify
  errors to prevent future recurrence. \emph{Example: training sessions and
  sharing of best practices.}\\[3pt]
  $\blacktriangleright$~\textbf{Directive controls:} provide guidance, processes,
  and training during duties. \emph{Example: detailed step-by-step instructions
  for specific processes.}};

\draw[-{Stealth[length=5pt]},forest,line width=1pt]
  ($(k1.north west)+(0.8,0)$) to[out=140,in=-70] ($(c2.south)+(0.6,0)$);
\draw[-{Stealth[length=5pt]},forest,line width=1pt]
  ($(k1.north east)+(-0.8,0)$) to[out=40,in=-110] ($(i2.south)+(-0.6,0)$);
\end{tikzpicture}
\end{center}
\orfigcap{Controls act on both the causes and the impacts of operational risk.}
